\appsection{Global Injectivity: Detailed Analysis}
\label{app:global_injectivity}
%==============================================================================
\label{sec:strategy}

For a scalar map $f \colon \R \to \R$, global injectivity is equivalent to
strict monotonicity, which can be checked from the sign of the derivative. For
maps $f \colon \R^n \to \R^n$ with $n > 1$, no such elementary criterion
exists: a nonvanishing Jacobian determinant guarantees only local injectivity,
and global conclusions require stronger conditions. The purpose of this
appendix, however, is narrower than identification in the strongest everywhere
sense: it is to provide sufficient conditions for the Jacobian-based
change-of-variables prior used in the Bayesian analysis. For that application,
injectivity is only needed on the support of the prior up to Lebesgue-null
exceptional sets, since such sets do not affect an absolutely continuous prior
density. The goal is therefore to show that the map from preference parameters
to certainty equivalents (CEs) and choice probabilities is injective on a
full-measure subset of the parameter space relevant for the transformed prior.
I begin with two general propositions that are used to provide the proof.

\begin{proposition}[Composition]
\label{prop:composition}
If $g \colon A \to B$ and $f \colon B \to C$ are both injective,
then the composite $f \circ g \colon A \to C$ is injective.
\end{proposition}

\begin{proof}
Suppose $f(g(y)) = f(g(y'))$. Injectivity of $f$ gives
$g(y) = g(y')$. Injectivity of $g$ gives $y = y'$.
\end{proof}

\begin{proposition}[Sequential decomposition]
\label{prop:sequential}
Let $f \colon X \times Y \to Z$, and for each $x \in X$ define the
conditional map $f_x \colon Y \to Z$ by $f_x(y) = f(x,y)$. Suppose:
\begin{enumerate}[(i)]
    \item For each $x \in X$, the conditional map $f_x$ is injective.
    Write $y = g(x,z)$ for the unique pre-image of each $z$ in the
    range of $f_x$.
    \item For each $z$ in the range of $f$, the reduced map
    $x \mapsto f(x,\, g(x,z))$ is injective.
\end{enumerate}
Then $f$ is injective.
\end{proposition}

\begin{proof}
Suppose $f(x,y) = f(x',y') = z$. Since $f_x(y) = z$, condition~(i)
gives $y = g(x,z)$; likewise $y' = g(x',z)$. Both $x$ and $x'$
satisfy $f(x,\, g(x,z)) = z$, so $x = x'$ by~(ii). Then
$y = g(x,z) = g(x',z) = y'$.
\end{proof}

The proposition reduces injectivity on a product space to two
lower-dimensional problems: identify~$y$ conditional on~$x$, then
identify~$x$ from the residual.


\subsection{Model, Conditions, Theorem, and Proof}
\label{sec:model}

The RDU model has three groups of parameters:
\begin{align}
    \bth &\in \Theta \subseteq \R^p
    &&\text{(utility parameters)}, \\
    \bomega &\in \mathcal{W} \subseteq \R^q
    &&\text{(PWF parameters)}, \\
    \lambda &\in (0, \infty)
    &&\text{(noise parameter)}.
\end{align}
The utility function $u(x; \bth)$ is strictly increasing and
differentiable in $x$ for all $\bth \in \Theta$. Let $a$ and $b$
denote the global minimum and maximum outcomes across all $K$
lotteries. Normalize utility on this common scale:
\begin{equation}\label{eq:normalization}
    \tilde{u}(x;\, \bth)
    = \frac{u(x;\, \bth) - u(a;\, \bth)}
           {u(b;\, \bth) - u(a;\, \bth)},
\end{equation}
so that $\tilde{u}(a;\, \bth) = 0$ and
$\tilde{u}(b;\, \bth) = 1$. Write
$\mu_j(\bth) = \tilde{u}(x_j;\, \bth) \in [0,1]$
for each outcome~$x_j$ appearing in any lottery.

The probability weighting function $w(p; \bomega)$ is strictly
increasing and differentiable in $p$ for all $\bomega \in \mathcal{W}$, with
$w(0; \bomega) = 0$ and $w(1; \bomega) = 1$. For each lottery $L_k$
with ranked outcomes and probabilities, the decumulative
probabilities and RDU decision weights are
\begin{align}
    S_i^{(k)} &= \sum_{j \geq i} p_j^{(k)}, \\
    \pi_i^{(k)}(\bomega) &= w(S_i^{(k)};\, \bomega)
      - w(S_{i+1}^{(k)};\, \bomega),
\end{align}
with the convention $w(0) = 0$ and $w(1) = 1$.

The normalized RDU value of lottery $k$ is
\begin{equation}
    V_k(\bth, \bomega)
    = \sum_{i=1}^{N_k} \pi_i^{(k)}(\bomega) \cdot \mu_i(\bth),
    \label{eq:rdu_value}
\end{equation}
where the sum runs over the outcomes of lottery~$k$, with
$\mu_a = 0$ and $\mu_b = 1$. The
certainty equivalent $c_k$ is the monetary amount satisfying
$\tilde{u}(c_k;\, \bth) = V_k$, i.e.,
$c_k = \tilde{u}^{-1}(V_k;\, \bth)$.

In a binary choice task, the subject observes lotteries $A$ and~$B$ and chooses
according to a noise model parameterized by~$\lambda$: the choice probability
$P = \Lambda\bigl((V_A - V_B)/(\lambda\,\Delta u)\bigr)$ depends on $(\bth,
\bomega)$ through the value difference and on~$\lambda$ through the noise
scaling, where $\Lambda$ is a strictly increasing link function and $\Delta u =
u(\xhi) - u(\xlo)$, where $\xhi$ and $\xlo$ are respectively the greatest and
least outcomes in lottery $k$.

The goal is to show that the extended map
\begin{equation}
    \widetilde{F} \colon (\bth, \bomega, \lambda) \;\mapsto\;
    (c_1, \ldots, c_K,\, P_{K+1})
\end{equation}
is injective on a full-measure subset of
$\Theta \times \mathcal{W} \times (0,\infty)$ sufficient for the
Bayesian change of variables, where $P_{K+1}$ is a choice probability
from an additional task.

For the Prelec--CRRA model of Section~\ref{sec:application_rdu}, three binary
lotteries and one choice probability suffice. Two binary lotteries with
distinct probabilities on the higher outcome produce two unique decision
weights, which is sufficient to identify the two Prelec parameters. With a
change of variables to decision weights, a third binary lottery that shares a
decision weight with $L_1$ or $L_2$ identifies the CRRA parameter. Because CRRA
is invariant to multiplicative scaling of outcomes, the outcomes of $L_3$
cannot be a multiplicative rescaling of the outcomes from the lottery with
which it shares a decision weight. The noise parameter is identified from
choice probabilities over any pair of the three lotteries. The remainder of
this section establishes sufficient conditions for injectivity on a
full-measure subset, which is the level needed for the transformed prior,
first generally for any utility function and probability weighting function
(Conditions~A, B, and~C below), followed by proofs that the proposed design satisfies them
(\cref{sec:condition_a_verification,sec:prelec_crra}) specifically for the CRRA
and 2 parameter Prelec functions.

\subsubsection{Lottery design}
\label{sec:design}

We use $K = p + q$ lotteries in two blocks:
$q$~Block~I lotteries $L_1, \ldots, L_q$ ($k = 1, \ldots, q$)
that identify the PWF parameters, and $p$~Block~II lotteries
($k = q{+}1, \ldots, p{+}q$) that identify the utility parameters.
Each lottery satisfies~\eqref{eq:rdu_value}; the partition into blocks
determines which parameters each group identifies.

The Block~I lotteries must collectively evaluate $w$ at $q$~distinct
probability points, enabling the factoring $\Phi_\mu = L \circ W$
(\cref{sec:same_cdf}). In the recommended \emph{all-binary design},
each Block~I lottery is binary (two-outcome) and evaluates $w$ at
exactly one probability point; with $q = 2$ binary lotteries having
distinct probabilities $p_1 \neq p_2$, the $L$-matrix is
diagonal and trivially invertible. Block~II provides additional
lotteries whose CEs depend on the utility parameters through a
channel not already captured by Block~I, identifying utility from
the residual equations. A single additional binary choice task
identifies the noise parameter~$\lambda$ (Condition~C).

\subsubsection{Condition~A: PWF identification from Block~I}
\label{sec:condition_a}

When $\bth$ is fixed, the normalized utilities
$\mu_j(\bth)$ are known constants, and the $q$~Block~I
instances of~\eqref{eq:rdu_value} become $q$~equations in
$q$~unknowns~$\bomega$. Condition~A requires that these
equations determine $\bomega$ uniquely---condition~(i) of
\cref{prop:sequential}.

\begin{condition}[A]
\label{cond:A}
For each fixed $\bth \in \Theta$, the map
$\Phi_\mu \colon \mathcal{W} \to \R^q$ defined by
\begin{equation}
    [\Phi_\mu(\bomega)]_k = \sum_i \pi_i^{(k)}(\bomega) \cdot \mu_i,
    \quad k = 1, \ldots, q
\end{equation}
is injective.
\end{condition}

When Condition~A holds, each $\bth$ determines a unique
$\bomega_A(\bth;\, \mathbf{c})$ solving the Block~I system.
The condition depends only on the PWF family and the Block~I
lottery structure, not on the observed CEs or the utility family.

\subsubsection{Condition~B: Utility identification from Block~II}
\label{sec:condition_b}

Substituting $\bomega = \bomega_A(\bth;\, \mathbf{c})$ from
Condition~A into the $p$~Block~II instances of~\eqref{eq:rdu_value}
eliminates the PWF parameters, yielding $p$~equations in
$\bth$ alone. Condition~B requires these residual equations
to determine $\bth$ uniquely---condition~(ii) of
\cref{prop:sequential}.

\begin{condition}[B]
\label{cond:B}
For any $\mathbf{c}$ in the range of~$F$, the reduced map
$H_{\mathbf{c}} \colon \Theta \to \R^p$ defined by
\begin{equation}
    [H_{\mathbf{c}}(\bth)]_k
    = \sum_i \pi_i^{(k)}\bigl(\bomega_A(\bth;\, \mathbf{c})\bigr)
      \cdot \mu_i(\bth)
    - \tilde{u}(c_k;\, \bth),
    \quad k = q+1, \ldots, p+q
\end{equation}
is injective.
\end{condition}

Unlike Condition~A, Condition~B depends on both the utility and PWF
families and on the observed CEs. When $\bth$ changes, three
things move simultaneously: the normalized utilities
$\mu_i(\bth)$, the implied PWF
$\bomega_A(\bth;\, \mathbf{c})$, and the left-hand side
$\tilde{u}(c_k;\, \bth)$. Condition~B requires that these
effects cannot cancel across all $p$ Block~II equations.

Condition~A implies that $W$ is injective on the Block~I
probability points (by contrapositive: identical decision weights
produce identical Block~I CEs). One may therefore reparameterize,
replacing $\bomega$ with the decision weights. This is most effective
when Block~II probabilities coincide with Block~I probabilities, so
that no inversion to PWF parameters is required.
\Cref{sec:condition_a_verification} establishes Condition~A for the
Prelec family, and \cref{sec:prelec_crra} exploits the resulting
$(r, w_1, w_2)$ reparameterization with a shared-probability
design ($p_3 = p_1$).

Verifying Condition~B requires the implicit derivative
$d\bomega_A/d\bth$, obtained by differentiating the Block~I
system:
\begin{equation}
    \frac{d\bomega_A}{d\bth}
    = -J_A^{-1}\,
    \frac{\partial \Phi_\mu}{\partial \bth},
    \label{eq:implicit_deriv}
\end{equation}
where $J_A = \partial \Phi_\mu / \partial \bomega$ (invertible by
Condition~A).

\subsubsection{Condition~C: Noise identification from a choice probability}
\label{sec:condition_c}

Once $(\bth, \bomega)$ are identified from the CEs, the value
difference $\Delta V = V_A - V_B$ for any additional choice task
is determined. Condition~C requires that the choice probability
$P = \Lambda(\Delta V / (\lambda\,\Delta u))$ identifies~$\lambda$.

\begin{condition}[C]
\label{cond:C}
For the identified $(\bth, \bomega)$, the choice task has
$\Delta V \neq 0$.
\end{condition}

\begin{proposition}[Noise identification]
\label{prop:noise}
If $\Lambda$ is strictly increasing and $\Delta V \neq 0$, then
$P = \Lambda(\Delta V / (\lambda\,\Delta u))$ is strictly monotone
in~$\lambda$, and Condition~C holds.
\end{proposition}

\begin{proof}
The argument $\Delta V / (\lambda\,\Delta u)$ is strictly monotone
in~$\lambda$ when $\Delta V \neq 0$; composing with the strictly
increasing~$\Lambda$ preserves strict monotonicity.
\end{proof}

Condition~C is pointwise, but the main text's
change-of-variables prior
$p(\bth,\bomega,\lambda)
= p_Y\!\bigl(\widetilde{F}(\cdot)\bigr)\,
\lvert\det J_{\widetilde{F}}\rvert$
requires injectivity to hold almost everywhere. Define the
indifference locus
$\mathcal{I} = \{\,(\bth, \bomega) : \Delta V(\bth, \bomega)
= 0\,\}$. For the prior density to be well defined,
$\mathcal{I}$ must have Lebesgue measure zero.
Degenerate designs---such as $A = B$, or pairs whose RDU values
coincide for all $(\bth, \bomega)$---give
$\mathcal{I} = \Theta \times \mathcal{W}$ and must be excluded. For
non-degenerate designs where $\Delta V$ is analytic and not
identically zero, $\mathcal{I}$ is a lower-dimensional analytic
variety (hence has Lebesgue measure zero), and Condition~C holds on
the full-measure complement. In practice, $A \neq B$ is sufficient
for standard PWF and utility families.

\subsubsection{Main theorem}
\label{sec:theorem}

\begin{theorem}[Injectivity sufficient for the Bayesian change of variables]
\label{thm:main}
Suppose Conditions~A and~B hold on $\Theta \times \mathcal{W}$, and let
\[
  \mathcal{I}
  = \{\,(\bth, \bomega) : \Delta V(\bth, \bomega) = 0\,\}.
\]
If $\mathcal{I}$ has Lebesgue measure zero, then the extended map
$\widetilde{F} \colon (\bth, \bomega, \lambda) \mapsto
(c_1, \ldots, c_{p+q},\, P_{K+1})$
is injective on the full-measure subset
\[
  \bigl((\Theta \times \mathcal{W}) \setminus \mathcal{I}\bigr)
  \times (0,\infty).
\]
Consequently, for any prior on $(\bth,\bomega,\lambda)$ that is absolutely
continuous with respect to Lebesgue measure, the transformed prior
\[
  p(\bth,\bomega,\lambda)
  = p_Y\!\bigl(\widetilde{F}(\bth,\bomega,\lambda)\bigr)\,
    \lvert \det J_{\widetilde{F}}(\bth,\bomega,\lambda)\rvert
\]
is well defined almost everywhere.
\end{theorem}

\begin{proof}
The CEs do not depend on~$\lambda$, so the map has
block-triangular structure: apply \cref{prop:sequential} with
$X = (\bth, \bomega)$, $Y = \lambda$, and
$Z = (c_1, \ldots, c_{p+q},\, P_{K+1})$.
Suppose
\[
  (\bth, \bomega),\; (\bth', \bomega') \notin \mathcal{I}
  \quad\text{and}\quad
  \widetilde{F}(\bth, \bomega, \lambda)
  = \widetilde{F}(\bth', \bomega', \lambda').
\]
By Condition~A, fixing $\bth$ determines
$\bomega = \bomega_A(\bth;\, \mathbf{c})$ uniquely, and likewise
$\bomega' = \bomega_A(\bth';\, \mathbf{c})$. Substituting into
the Block~II equations yields $p$~equations in $\bth$ alone;
by Condition~B, $\bth = \bth'$, and then $\bomega = \bomega'$.
With $(\bth, \bomega)$ identified, $\Delta V$ is determined. Because
$(\bth,\bomega) \notin \mathcal{I}$, we have $\Delta V \neq 0$, so
\cref{prop:noise} gives $\lambda = \lambda'$.
Thus $\widetilde{F}$ is injective on
$((\Theta \times \mathcal{W}) \setminus \mathcal{I}) \times (0,\infty)$.
The final statement follows because an absolutely continuous prior assigns
probability zero to $\mathcal{I} \times (0,\infty)$, so null exceptional sets
do not affect the induced prior measure.
\end{proof}


\subsection{Verifying Condition~A: Two-Point Injectivity and the $L \circ W$ Factoring}
\label{sec:condition_a_verification}

Conditions~A and~B require verification for specific utility and
PWF families; Condition~C requires a non-degenerate choice task
($\Delta V \neq 0$). This section and the next verify all three
conditions for Prelec--CRRA.

For $q = 2$, the map $\Phi_\mu$ factors as $L \circ W$, where
$W$ is the two-point evaluation map and $L$ is a $2 \times 2$
matrix determined by the lottery structure. When $\det(L) \neq 0$,
Condition~A reduces to injectivity of~$W$.

\subsubsection{Two-point injectivity}
\label{sec:two_point}

For the Prelec PWF
$w(p;\, \phi, \eta) = \exp\bigl(-\eta\, t^\phi\bigr)$ with
$t = -\ln p > 0$, define the \emph{two-point evaluation map}
$W \colon (\phi, \eta) \mapsto \bigl(w(p_1),\, w(p_2)\bigr)$
for distinct $p_1 \neq p_2 \in (0,1)$, and write
$t_i = -\ln p_i$, $w_i = w(p_i)$.

\begin{lemma}[Prelec two-point injectivity]
\label{lem:prelec_two_point}
The two-point evaluation map $W$ is
injective on $\R^2_{>0}$ for any $p_1 \neq p_2 \in (0,1)$.
\end{lemma}

\begin{proof}
Write $t_i = -\ln p_i > 0$. Suppose
$W(\phi_1, \eta_1) = W(\phi_2, \eta_2)$. Taking logarithms of both
components:
\begin{align}
    \eta_1\, t_1^{\phi_1} &= \eta_2\, t_1^{\phi_2},
    \label{eq:prelec_eq1} \\
    \eta_1\, t_2^{\phi_1} &= \eta_2\, t_2^{\phi_2}.
    \label{eq:prelec_eq2}
\end{align}
Dividing~\eqref{eq:prelec_eq1} by~\eqref{eq:prelec_eq2}:
\begin{equation}
    \left(\frac{t_1}{t_2}\right)^{\!\phi_1}
    = \left(\frac{t_1}{t_2}\right)^{\!\phi_2}.
    \label{eq:ratio}
\end{equation}
Since $p_1 \neq p_2$, the values $t_1 \neq t_2$, so
$t_1/t_2 \neq 1$. The function $a^x$ is injective in~$x$ for any
$a > 0$ with $a \neq 1$, so~\eqref{eq:ratio} implies
$\phi_1 = \phi_2$. Substituting back
into~\eqref{eq:prelec_eq1}: $\eta_1 = \eta_2$, since
$t_1^{\phi_1} > 0$.
\end{proof}

The proof uses only the \emph{ratio property}: a logarithmic
transformation separates the parameters multiplicatively, and the
resulting basis functions form a ratio differing from~$1$ at
distinct probability points. Any two-parameter PWF with this
structure (e.g., Goldstein--Einhorn) satisfies two-point
injectivity by the same argument.


\subsubsection{$W$ is a bijection: the $(r, w_1, w_2)$ reparameterization}
\label{sec:W_bijection}

Injectivity of~$W$ alone enables the change of variables from
$(\phi, \eta)$ to $(w_1, w_2)$: one works in the image
$W(\R^2_{>0}) \subseteq (0,1)^2$, which may be a proper subset,
and every statement about $(r, w_1, w_2)$ applies to the parameter
values actually attained by the PWF family. The simplification
from working in decision-weight space is substantial enough to
justify this reparameterization even when the image $W(\R^2_{>0})$
must be characterized explicitly. When $W$ is bijective, the image
equals the full codomain and no such characterization is needed.

For the Prelec family, $W$ is in fact \emph{bijective} onto the
order-compatible region
\begin{equation}
  W(\R^2_{>0}) =
  \begin{cases}
    \{(w_1, w_2) : 0 < w_1 < w_2 < 1\}, & \text{if } p_1 < p_2, \\[4pt]
    \{(w_1, w_2) : 0 < w_2 < w_1 < 1\}, & \text{if } p_1 > p_2.
  \end{cases}
\end{equation}
The proof of \cref{lem:prelec_two_point} provides the explicit
inverse. Given observed values $w_1 = w(p_1)$ and
$w_2 = w(p_2)$ with $t_i = -\ln p_i$:
\begin{equation}\label{eq:W_inverse}
    \phi = \frac{\ln(\ln w_1 / \ln w_2)}{\ln(t_1 / t_2)},
    \qquad
    \eta = \frac{-\ln w_1}{t_1^{\,\phi}}.
\end{equation}
Both expressions are well defined whenever $w_1, w_2 \in (0,1)$
satisfy the ordering induced by $p_1$ and $p_2$. The condition
$p_1 \neq p_2$ ensures $t_1 \neq t_2$ so the denominator is nonzero,
and $w_i \in (0,1)$ ensures $-\ln w_i > 0$. If $p_1 < p_2$, then
$t_1 > t_2$, so $\ln(t_1/t_2) > 0$ and positivity of $\phi$ requires
$\ln(\ln w_1 / \ln w_2) > 0$, equivalently $w_1 < w_2$.
The case $p_1 > p_2$ is symmetric, yielding $w_1 > w_2$.
This bijectivity is likely available for other common two-parameter
PWF families (e.g., Goldstein--Einhorn). Because of this bijection,
we conduct all subsequent analysis in $(r, w_1, w_2)$ space subject
only to the ordering restriction implied by $p_1$ and $p_2$; the
Prelec parameters do not appear again.

For the all-binary design, each Block~I lottery $L_k$ is binary with
probability $p_k$ on the better outcome. When $p_1 \neq p_2$,
the $L$-matrix in the $L \circ W$ factoring is
$L = \diag(\Delta\mu_1, \Delta\mu_2)$ with
$\det(L) = \Delta\mu_1\,\Delta\mu_2 > 0$
(\cref{sec:same_cdf} below), and Condition~A holds
unconditionally.

\begin{proposition}[All-binary Condition~A]
\label{prop:binary}
For the all-binary Block~I design with $p_1 \neq p_2$,
Condition~A holds at all $(\phi, \eta) \in \R^2_{>0}$ and all
utility parameters, with no restrictions on outcomes.
\end{proposition}

\begin{proof}
Each binary lottery evaluates $w$ at one probability, giving
$L = \diag(\Delta\mu_1, \Delta\mu_2)$ with $\det(L) > 0$.
By \cref{lem:prelec_two_point} and the invertibility of $L$,
Condition~A holds.
\end{proof}

\subsubsection{The $L \circ W$ factoring}
\label{sec:same_cdf}

For a binary lottery $L_k$ with outcomes $(\xlo_k, \xhi_k)$
and probability $p_k$ on the better outcome, the normalized RDU
value~\eqref{eq:rdu_value} is
\begin{equation}\label{eq:binary_rdu}
    V_k = \mu_{\xlo_k} + w(p_k)\,\Delta\mu_k,
    \qquad
    \Delta\mu_k \coloneqq \mu_{\xhi_k} - \mu_{\xlo_k} > 0.
\end{equation}
With two binary lotteries having $p_1 \neq p_2$:
\begin{equation}
    \begin{pmatrix} V_1 - \mu_{\xlo_1} \\ V_2 - \mu_{\xlo_2} \end{pmatrix}
    = \underbrace{
      \begin{pmatrix} \Delta\mu_1 & 0 \\ 0 & \Delta\mu_2 \end{pmatrix}
    }_{L}
    \underbrace{
      \begin{pmatrix} w(p_1) \\ w(p_2) \end{pmatrix}
    }_{W(\phi, \eta)},
    \label{eq:factoring_binary}
\end{equation}
with $\det(L) = \Delta\mu_1\,\Delta\mu_2 > 0$. For fixed
$\bth$, Condition~A reduces to injectivity of $W$, which holds
by \cref{lem:prelec_two_point}.


\subsection{Application: Prelec--CRRA in $(r, w_1, w_2)$ Space}
\label{sec:prelec_crra}

For the concrete case in the main text:
$\bth = r$ (CRRA coefficient, $p = 1$) and
$\bomega = (\phi, \eta)$ (Prelec parameters, $q = 2$), with utility
$u(x; r) = x^{1-r}/(1-r)$ and PWF
$w(p; \phi, \eta) = \exp\bigl(-\eta(-\ln p)^\phi\bigr)$.
By the $W$-bijectivity of \cref{sec:W_bijection}, the model is
reparameterized as $F \colon (r, w_1, w_2) \mapsto (c_1, c_2, c_3)$.
All analysis in this section operates in $(r, w_1, w_2)$ space.


\subsubsection{The CE map}
\label{sec:ce_function}

For a binary lottery $L_k = (\xlo_k, \xhi_k)$ with decision
weight~$w_k$ on the better outcome, write $s = 1 - r$. The CRRA
certainty equivalent is a power mean:
\begin{equation}
  \label{eq:ce_decomp}
  c_k^{\,s} = w_k\, \xhi_k^s + (1 - w_k)\, \xlo_k^s,
\end{equation}
equivalently $c_k = [w_k\, \xhi_k^s + (1-w_k)\, \xlo_k^s]^{1/s}$.
Since the right-hand side is strictly increasing in~$w_k$,
Condition~A is immediate in $(r, w_1, w_2)$ space: given~$r$,
each $c_k$ is uniquely determined by~$w_k$.

\subsubsection{Condition~B: CRRA identification via shared probability}
\label{sec:cond_b_crra}

With $p = 1$, Condition~B reduces to the scalar condition
$dc_3/dr\big|_{c_1} \neq 0$. The shared-probability design
$p_3 = p_1$ eliminates the PWF parameters entirely:
both $L_1$ and $L_3$ share the decision weight $w_3 = w_1$.

Parameterize $L_3$'s outcomes relative to~$L_1$'s:
\begin{equation}\label{eq:affine}
    \xhi_3 = m\,\xhi_1 + b, \qquad \xlo_3 = m\,\xlo_1 + b,
\end{equation}
where $m = (\xhi_3 - \xlo_3)/(\xhi_1 - \xlo_1) > 0$ is the spread ratio
and $b = \xlo_3 - m\,\xlo_1$ is the shift.
Substituting into~\eqref{eq:ce_decomp} with $w_3 = w_1$:
\begin{align}
    c_1^{\,s} &= w_1\, \xhi_1^s + (1-w_1)\, \xlo_1^s,
    \label{eq:c1_shared} \\
    c_3^{\,s} &= w_1\, (m\,\xhi_1 + b)^s + (1-w_1)\, (m\,\xlo_1 + b)^s.
    \label{eq:c3_shared}
\end{align}
Solving~(\ref{eq:c1_shared}) for $w_1$ in terms of $s$:
\begin{align}
  \label{eq:solved_w1}
  w_1(s) &= \frac{c_1^{\,s} - \xlo_1^s}{\xhi_1^s - \xlo_1^s}.
\end{align}
Substituting into~(\ref{eq:c3_shared}) makes $c_3$ a function of $s$ alone:
\begin{align}
  \label{eq:c3_wr}
  c_3^{\,s} &= w_1(s) \left[ (m\,\xhi_1 + b)^s - (m\,\xlo_1 + b)^s\right] + (m\,\xlo_1 + b)^s.
\end{align}

Setting $b = 0$ reveals a necessary condition. Expanding $w_1(s)$ and simplifying:
\begin{align}
  c_3^{\,s} &= w_1(s) \left[ m^s\,\xhi_1^s - m^s\,\xlo_1^s\right] + m^s\,\xlo_1^s \notag \\
  c_3^{\,s} &= w_1(s) m^s\left[ \xhi_1^s - \xlo_1^s\right] + m^s\,\xlo_1^s \notag \\
  c_3^{\,s} &= m^s \left[w_1(s) \left( \xhi_1^s - \xlo_1^s\right) + \xlo_1^s\right] \notag \\
  c_3^{\,s} &= m^s \left[\frac{c_1^{\,s} - \xlo_1^s}{\, \xhi_1^s - \xlo_1^s} \left( \xhi_1^s - \xlo_1^s\right) + \xlo_1^s\right] \notag \\
  c_3^{\,s} &= m^s \left[c_1^{\,s} - \xlo_1^s + \xlo_1^s\right] \notag \\
  c_3^{\,s} &= (m\, c_1)^{\,s} \notag \\
  c_3 &= m\, c_1
\end{align}
When $b = 0$, the CE $c_3 = m\,c_1$ is independent of~$s$, so $L_3$ provides
no information about~$r$. This reflects the well-known scale invariance of
CRRA utility. The remainder of this section shows that $b \neq 0$ is also
\textit{sufficient} for identification.

CRRA homogeneity $\CE(mL) = m \cdot \CE(L)$ implies that the scale
factor~$m$ does not affect monotonicity of $c_3$ in~$r$, so we may
normalize by dividing $L_3$'s outcomes by~$m$. Writing $e = b/m$, the
condition $b \neq 0$ is equivalent to $e \neq 0$, and the CE equations become
\begin{align}
    c_1^{\,s} &= w_1\, \xhi_1^s + (1-w_1)\, \xlo_1^s,
    \label{eq:c1_shift} \\
    \tilde{c}_3^{\,s} &= w_1\, (\xhi_1 + e)^s
      + (1-w_1)\, (\xlo_1 + e)^s,
    \label{eq:c3_shift}
\end{align}
where $\tilde{c}_3$ is the CE of the normalized lottery (with
$c_3 = m\,\tilde{c}_3$); write $c_3$ for $\tilde{c}_3$ hereafter.

Implicitly differentiating~(\ref{eq:c3_shift}) with respect to~$s$:
\begin{align}
  \label{eq:idiff}
  \begin{split}
  c_3^{\,s}\!\left(\ln c_3 + \frac{s}{c_3}\frac{dc_3}{ds}\right) &=
    w_1'(s)\,[(\xhi_1{+}e)^s - (\xlo_1{+}e)^s]
    + w_1(s)\,[(\xhi_1{+}e)^s \ln(\xhi_1{+}e) - (\xlo_1{+}e)^s \ln(\xlo_1{+}e)] \\
    &\quad + (\xlo_1{+}e)^s \ln(\xlo_1{+}e)
  \end{split}
\end{align}
From~(\ref{eq:solved_w1}), the quotient rule gives:
\begin{align}
  w_1'(s) &= \frac{(c_1^{\,s}\!\ln c_1 - \xlo_1^s\!\ln\xlo_1)(\xhi_1^s - \xlo_1^s)
    - (c_1^{\,s} - \xlo_1^s)(\xhi_1^s\!\ln\xhi_1 - \xlo_1^s\!\ln\xlo_1)}
    {(\xhi_1^s - \xlo_1^s)^2} \notag \\[0.5em]
  \label{eq:dw1ds}
  -w_1'(s) &= \frac{w_1\,\xhi_1^s \ln \xhi_1
    + (1-w_1)\,\xlo_1^s \ln \xlo_1
    - c_1^s \ln c_1}{\xhi_1^s - \xlo_1^s}.
\end{align}
Solving~(\ref{eq:idiff}) for $dc_3/ds$:
\begin{align}
  \label{eq:dc3ds}
  \frac{dc_3}{ds} = \frac{1}{s\, c_3^{\,s-1}}\Bigl[&\;
    w_1'(s)\,[(\xhi_1{+}e)^s - (\xlo_1{+}e)^s] \notag \\
    &\; + w_1(s)\,[(\xhi_1{+}e)^s \ln(\xhi_1{+}e)
      - (\xlo_1{+}e)^s \ln(\xlo_1{+}e)]
    + (\xlo_1{+}e)^s \ln(\xlo_1{+}e) - c_3^{\,s} \ln c_3
  \Bigr]
\end{align}
Condition~B holds if $dc_3/ds \neq 0$ for all $s \neq 0$
(equivalently $dc_3/dr \neq 0$, since $s = 1 - r$).
To analyze the sign of the bracket in~(\ref{eq:dc3ds}),
define the Jensen gap of $\phi(t) = t\ln t$ for lottery~$L_k$:
\begin{equation}\label{eq:jensen_gap_def}
  G_k
  = w_1\,\phi(\xhi_k^s) + (1-w_1)\,\phi(\xlo_k^s)
    - \phi(c_k^s),
\end{equation}
where $c_k^s = w_1\,\xhi_k^s + (1-w_1)\,\xlo_k^s$.
Since $\phi$ is strictly convex, $G_k > 0$ for all $s$.
Using $\phi(x^s) = x^s \ln(x^s) = s\, x^s \ln x$,
the three non-$w_1'$ terms in the bracket
of~(\ref{eq:dc3ds}) are
\begin{equation}\label{eq:bracket_J3}
  w_1\,(\xhi_1{+}e)^s \ln(\xhi_1{+}e)
  + (1-w_1)\,(\xlo_1{+}e)^s \ln(\xlo_1{+}e)
  - c_3^s \ln c_3
  \;=\; \frac{G_3}{s}.
\end{equation}
Similarly, from~(\ref{eq:dw1ds}), the numerator of
$-w_1'(s)$ is $G_1/s$ and the denominator is
$\xhi_1^s - \xlo_1^s$, so
\begin{equation}\label{eq:w1prime_J1}
  -w_1'(s) = \frac{G_1}
    {s\,(\xhi_1^s - \xlo_1^s)}.
\end{equation}
Substituting into~(\ref{eq:dc3ds}):
\begin{equation}\label{eq:dc3ds_collapsed}
  \frac{dc_3}{ds}
  = \frac{1}{s\, c_3^{s-1}}
    \left[
      \frac{-G_1\bigl[(\xhi_1{+}e)^s - (\xlo_1{+}e)^s\bigr]}
           {s\,(\xhi_1^s - \xlo_1^s)}
    + \frac{G_3}{s}\right].
\end{equation}
Since $s^2 > 0$ and $c_3^{s-1} > 0$, the sign of
$dc_3/ds$ is determined by the bracket, so
Condition~B is equivalent to
\begin{equation}\label{eq:condB_ratio}
  \frac{w_1\,\xhi_1^s\ln\xhi_1^s
    + (1{-}w_1)\,\xlo_1^s\ln\xlo_1^s
    - c_1^s\ln c_1^s}
       {\xhi_1^s - \xlo_1^s}
  \;\neq\;
  \frac{w_1\,(\xhi_1{+}e)^s\ln(\xhi_1{+}e)^s
    + (1{-}w_1)\,(\xlo_1{+}e)^s\ln(\xlo_1{+}e)^s
    - c_3^s\ln c_3^s}
       {(\xhi_1{+}e)^s - (\xlo_1{+}e)^s}
\end{equation}
for all $s \neq 0$.
Suppose for contradiction that equality holds
in~(\ref{eq:condB_ratio}) for some $s \neq 0$ and
$e \neq 0$.
Using $x^s\ln(x^s) = s\,x^s\ln x$, a common factor
of~$s$ appears in both numerators and cancels, giving
\begin{equation}\label{eq:condB_s_cancelled}
  \frac{w_1\,\xhi_1^s\ln\xhi_1
    + (1{-}w_1)\,\xlo_1^s\ln\xlo_1
    - c_1^s\ln c_1}
       {\xhi_1^s - \xlo_1^s}
  \;=\;
  \frac{w_1\,(\xhi_1{+}e)^s\ln(\xhi_1{+}e)
    + (1{-}w_1)\,(\xlo_1{+}e)^s\ln(\xlo_1{+}e)
    - c_3^s\ln c_3}
       {(\xhi_1{+}e)^s - (\xlo_1{+}e)^s}.
\end{equation}
Write $\rho_k = (\xlo_k/\xhi_k)^s$ and divide
numerator and denominator on each side
of~(\ref{eq:condB_s_cancelled}) by~$\xhi_k^s$. Since
$\ln\xlo_k = \ln\xhi_k + \tfrac{1}{s}\ln\rho_k$
and $\ln c_k = \ln\xhi_k
  + \tfrac{1}{s}\ln[w_1{+}(1{-}w_1)\rho_k]$,
the coefficient of~$\ln\xhi_k$ in each numerator is
$w_1 + (1{-}w_1)\rho_k - [w_1{+}(1{-}w_1)\rho_k] = 0$.
A common factor of~$\tfrac{1}{s}$ cancels, reducing~(\ref{eq:condB_s_cancelled}) to
\begin{equation}\label{eq:ratio_of_rho}
  \frac{(1{-}w_1)\,\rho_1\ln\rho_1
    - [w_1 {+} (1{-}w_1)\rho_1]\,
      \ln[w_1 {+} (1{-}w_1)\rho_1]}
       {1 - \rho_1}
  \;=\;
  \frac{(1{-}w_1)\,\rho_3\ln\rho_3
    - [w_1 {+} (1{-}w_1)\rho_3]\,
      \ln[w_1 {+} (1{-}w_1)\rho_3]}
       {1 - \rho_3}.
\end{equation}
Both sides are the same function of~$\rho_k$
(with~$w_1$ fixed), evaluated at $\rho_1$
and~$\rho_3$ respectively. We show this function is
strictly decreasing, so the equality forces
$\rho_1 = \rho_3$.

Write $f(\rho)$ for the numerator
of~(\ref{eq:ratio_of_rho}), so the ratio is
$f(\rho)/(1-\rho)$. Differentiating $f$ term by
term:
\begin{align}
  \frac{d}{d\rho}\bigl[(1{-}w_1)\,\rho\ln\rho\bigr]
    &= (1{-}w_1)(\ln\rho + 1),
  \notag\\[4pt]
  \frac{d}{d\rho}
  \bigl\{
    {-}[w_1{+}(1{-}w_1)\rho]\,
    \ln[w_1{+}(1{-}w_1)\rho]
  \bigr\}
    &= -(1{-}w_1)
       \bigl(\ln[w_1{+}(1{-}w_1)\rho] + 1\bigr).
  \notag
\end{align}
The constant terms $\pm(1{-}w_1)$ cancel, leaving
\begin{equation}\label{eq:f_prime}
  f'(\rho)
  = (1{-}w_1)\,
    \ln\frac{\rho}{w_1 + (1{-}w_1)\rho}.
\end{equation}
Since $\rho < w_1 + (1{-}w_1)\rho$ for all
$w_1 \in (0,1)$ and $\rho > 0$, the logarithm is
negative, so $f'(\rho) < 0$ everywhere.

By the quotient rule, the derivative of
$f(\rho)/(1-\rho)$ has the same sign as
\begin{equation}\label{eq:qr_num}
  f'(\rho)\,(1 - \rho) \;+\; f(\rho),
\end{equation}
since the denominator squared $(1-\rho)^2$ is
positive. We claim~(\ref{eq:qr_num}) is negative for
all $\rho \neq 1$.

At $\rho = 1$, $f(1) = 0$ since both terms vanish
($\ln 1 = 0$). Combined with $(1 - \rho) = 0$,
expression~(\ref{eq:qr_num}) vanishes at $\rho = 1$.

Differentiating~(\ref{eq:qr_num}):
\begin{equation}\label{eq:qr_deriv}
  \frac{d}{d\rho}\bigl[f'(\rho)(1{-}\rho) + f(\rho)\bigr]
  = f''(\rho)(1{-}\rho)
    \underbrace{{}-f'(\rho) + f'(\rho)}_{=\,0}
  = f''(\rho)(1{-}\rho).
\end{equation}
From~(\ref{eq:f_prime}),
\begin{equation}\label{eq:f_double_prime}
  f''(\rho)
  = (1{-}w_1)\biggl[
      \frac{1}{\rho}
      - \frac{1{-}w_1}{w_1{+}(1{-}w_1)\rho}
    \biggr]
  = \frac{w_1(1{-}w_1)}
         {\rho\,[w_1{+}(1{-}w_1)\rho]}
  \;>\; 0.
\end{equation}
So by~(\ref{eq:qr_deriv}), the derivative
of~(\ref{eq:qr_num}) is $f''(\rho)(1-\rho)$: positive
for $\rho < 1$, negative for $\rho > 1$.
Expression~(\ref{eq:qr_num}) therefore has its unique
maximum at $\rho = 1$ where it equals zero, hence is
strictly negative for all $\rho \neq 1$, and
$f(\rho)/(1-\rho)$ is strictly decreasing on
$(0,\infty) \setminus \{1\}$.

\begin{proposition}[Shared-probability Condition~B for CRRA]
\label{prop:shared_prob_condB}
Under the shared-probability design $p_3 = p_1$, if
$\xhi_1/\xlo_1 \neq \xhi_3/\xlo_3$ then Condition~B holds for all
$s \neq 0$, equivalently for all $r \neq 1$.
\end{proposition}

\begin{proof}
The equality in~(\ref{eq:ratio_of_rho}) forces
$\rho_1 = \rho_3$, i.e.,
\begin{equation}\label{eq:ratio_equality}
  \biggl(\frac{\xlo_1}{\xhi_1}\biggr)^{\!s}
  =
  \biggl(\frac{\xlo_1 + e}{\xhi_1 + e}\biggr)^{\!s}.
\end{equation}
For $s \neq 0$,~(\ref{eq:ratio_equality}) gives
$\xlo_1/\xhi_1 = (\xlo_1{+}e)/(\xhi_1{+}e)$.
Cross-multiplying:
$\xlo_1(\xhi_1{+}e) = \xhi_1(\xlo_1{+}e)$,
which simplifies to $e\,(\xhi_1 - \xlo_1) = 0$.
Since $\xhi_1 \neq \xlo_1$, this forces $e = 0$,
contradicting our assumption. Thus Condition~B holds for all $s \neq 0$.
\end{proof}


\subsubsection{Full three-parameter identification}
\label{sec:full_identification}

Combining the results: the $W$-bijectivity of
\cref{sec:W_bijection} and the diagonal $L$-matrix establish
Condition~A (\cref{prop:binary}) for any $p_1 \neq p_2$.
The shared-probability design $p_3 = p_1$ with
$\xhi_1/\xlo_1 \neq \xhi_3/\xlo_3$ establishes Condition~B
(\cref{prop:shared_prob_condB}) on the power-CRRA domain
$\R_{>0} \setminus \{1\}$, equivalently $s = 1-r \neq 0$.
Thus the CE map
$F \colon (r, w_1, w_2) \mapsto (c_1, c_2, c_3)$
is injective on the full-measure subset
\[
  \bigl((0,\infty) \setminus \{1\}\bigr) \times W(\R^2_{>0}).
\]
The excluded singleton $r = 1$ corresponds to the log-utility limit and is
immaterial for the Bayesian change-of-variables prior.


\subsubsection{Condition~C: noise identification}
\label{sec:noise}

The main text uses a logistic Fechner model
with $\Lambda(t) = 1/(1 + e^{-t})$. Condition~C
(\cref{sec:condition_c}) holds for any choice task with
$\Delta V \neq 0$. For a non-degenerate design, the indifference locus
$\Delta V = 0$ is a lower-dimensional exceptional set, so combined with
Conditions~A and~B the full map
$\widetilde{F} \colon (r, w_1, w_2, \lambda) \mapsto
(c_1, c_2, c_3,\, P_4)$ is injective on the full-measure subset relevant for
Bayesian estimation. This is exactly the level of injectivity needed to justify
the transformed prior in the main text.


\subsection{Extensions and Discussion}
\label{sec:inj_extensions}

The sequential approach scales linearly: $q$ Block~I lotteries plus
$p$ Block~II lotteries, plus one choice task for~$\lambda$.
The all-binary design handles Block~I for
$q = 2$ (covering Prelec, Goldstein--Einhorn, Tversky--Kahneman);
Block~II lotteries are chosen freely to provide outcome-side
variation. Condition~C requires one additional choice task with
$\Delta V \neq 0$.
For utility families with $p > 1$ (e.g., expo-power),
Condition~B becomes a $p \times p$ injectivity condition.

The change-of-variables computation in the main text requires the
Jacobian of the Condition~A map, which factors as
$J_A = L \cdot J_W$ by the chain rule. That computation uses the
original $(\phi, \eta)$ parameterization and is carried out in
Appendix~\ref{app:jacobian_details}.



\pagebreak
